% Gauss, Werke X/1 (1917), native123–126, printed116–119.
\SourceLeaf{123}\NoSourceHead
\SourceGroup{III}
\HistoricTitle{[DIE UNZERLEGBARKEIT DER\\KREISTHEILUNGSGLEICHUNG.]}
\HistoricCitation{[Aus Handbuch 18, Bd, October 1805, S. 198–199.]}
\NumberedItem{1}
\noindent\hspace*{\parindent}Die radices propriae der Gleichung
\[x^{a^\alpha b^\beta c^\gamma d^\delta\cdots}-1=0\]
\noindent oder
\[x^n-1=0,\]
\noindent wo $a,b,c,d$ etc. ungleiche Primzahlen bedeuten, sind enthalten in der Gleichung
\begin{center}
\resizebox{\linewidth}{!}{$\displaystyle
\frac{x^n-1\mathpunct{.}\,x^{\frac n{ab}}-1\mathpunct{.}\,x^{\frac n{ac}}-1\mathpunct{.}\,x^{\frac n{bc}}-1\text{ etc. }x^{\frac n{abcd}}-1\text{ etc. etc.}}
{x^{\frac na}-1\mathpunct{.}\,x^{\frac nb}-1\mathpunct{.}\,x^{\frac nc}-1\mathpunct{.}\,x^{\frac nd}-1\text{ etc. }x^{\frac n{abc}}-1\mathpunct{.}\,x^{\frac n{abd}}-1\mathpunct{.}\,x^{\frac n{bcd}}-1\text{ etc. etc.}}=0$}
\end{center}
\noindent oder
\[X=0.\]
\noindent Der Werth von $X$ für $x=1$ wird gleich 1, wenn mehr als ein Factor im Exponenten
$n$, oder $=a$, wenn $n=a^\alpha$. Setzt man die Anzahl der Factoren $a,b,c,d$ etc. $=\nu$,
so sind
\[
\left.\begin{array}{r}\text{im Zähler}\\[14pt]\text{im Nenner}\end{array}\right\}
\text{von }X
\left\{\begin{array}{l}
1+\dfrac{\nu\mathpunct{.}\nu-1}{1\mathpunct{.}2}+\dfrac{\nu\mathpunct{.}\nu-1\mathpunct{.}\nu-2\mathpunct{.}\nu-3}{1\mathpunct{.}2\mathpunct{.}3\mathpunct{.}4}+\text{etc.}\\[14pt]
\nu+\dfrac{\nu\mathpunct{.}\nu-1\mathpunct{.}\nu-2}{1\mathpunct{.}2\mathpunct{.}3}+\ldots
\end{array}\right.
\]
\noindent Factoren, oder in jedem $2^{\nu-1}$.
\clearpage
\SourceLeaf{124}\ArithmeticHead{117}{R}{DIE UNZERLEGBARKEIT DER KREISTHEILUNGSGLEICHUNG.}
\noindent\hspace*{\parindent}Wenn $X$ einen Factor vom Grade $\lambda$ hätte, dessen Coefficienten
alle rational wären, so lässt sich zeigen

1) dass $\lambda$ durch $a^{\alpha-1}(a-1)$, $b^{\beta-1}(b-1)$, $c^{\gamma-1}(c-1)$
u. s. w. theilbar seyn müsste,

2) dass sodann auch die Function, deren Wurzeln alle eigentlichen Wurzeln der Gleichung
\[x^{b^\beta c^\gamma d^\delta\cdots}-1=0\]
\noindent sind, einen Factor der Dimension $\dfrac\lambda{a^{\alpha-1}(a-1)}$ haben müsste,
dessen Coefficienten alle durch Wurzeln der Gleichung
\[x^{a^\alpha}-1=0\]
\noindent rational dargestellt werden könnten.

Von dem Beweise würde folgendes Hauptmomente ausmachen.

I. Wenn $a,b,c,d$ etc. ganze Zahlen sind, $R$ Wurzel der Gleichung $x^\beta-1=0$
und $p$ Primzahl $=k\beta+1$, so wird
\[(a+bR+cR^2+dR^3+\ldots)^\beta\]
\noindent entwickelt nicht
\[\equiv0\ (\text{mod. }p,\ x^{\beta-1},\ x^{\beta-2},\ldots)\]
\noindent seyn können.

\textls[120]{Bew}[eis]: Wäre
\[(a+bR+cR^2+\ldots)^\beta\equiv0\]
\noindent so wäre auch
\[(a+bR+cR^2+\ldots)^{\beta k+1}\equiv0;\]
\noindent aber jene Potenz ist
\[\equiv a+bR+cR^2+\text{etc.}\equiv(b-a)R+(c-a)R^2+\text{etc.},\]
\noindent folglich müsste sein
\[a\equiv b\equiv c\text{ etc.},\]
\noindent daraus würde aber folgen
\[a^2+b^2+c^2+\ldots-ab-ac-\text{etc.}\equiv0\ (\text{mod. }pp).\]
\clearpage
\SourceLeaf{125}\ArithmeticHead{118}{L}{ALGEBRA.\quad NACHLASS.}
\NumberedItem{2}
\noindent\hspace*{\parindent}Gleichung
\[x^{a^\alpha{a'}^{\alpha'}{a''}^{\alpha''}{a'''}^{\alpha'''}}-1=0\]
\noindent oder
\[x^n-1=0.\]
\noindent [Die] Gleichung der primitiven Wurzeln $X=0$ hat
\[n\frac{a-1\mathpunct{.}a'-1\mathpunct{.}a''-1\ldots}{a\mathpunct{.}a'\mathpunct{.}a''\ldots}\]
\noindent Dimensionen. $X$ habe den rationalen Divisor $Y$, und es sei [$\rho$] eine Wurzel
der Gleichung $Y=0$, [für] diese $\rho$ kann man voraussetzen
\[\rho=rr'r''\ldots,\]
\noindent so dass $r,r',r''$ u. s. w. primitive Wurzeln der Gleichungen
\[x^{a^\alpha}-1=0,\quad x^{{a'}^{\alpha'}}-1=0\text{ etc.}\]
\noindent sind. Es sei ferner $\sigma$ eine Wurzel der Gleichung
\[\frac XY=0,\]
\noindent so kann man voraussetzen
\[\sigma=r^\beta{r'}^{\beta'}{r''}^{\beta''}\ldots,\]
\noindent so dass $\beta,\beta',\beta''$ etc. Divisoren von
\[a^{\alpha-1}(a-1),\quad{a'}^{\alpha'-1}(a'-1),\quad{a''}^{\alpha''-1}(a''-1)\text{ etc.}\]
\noindent sind.

Alle Wurzeln, die ein vollständiges System bilden, werden sich so darstellen lassen:
\[g,g',g'',\ldots\text{ radices primit. congr. }x^{a^{\alpha-1}(a-1)}\equiv1\ (\text{mod. }a^\alpha)\text{ etc.}\]
\[r^{g^{\lambda x}}{r'}^{{g'}^{\mu x+\mu'x'}}{r''}^{{g''}^{\nu x+\nu'x'+\nu''x''}}\ldots,\]
\clearpage
\SourceLeaf{126}\ArithmeticHead{119}{R}{DIE UNZERLEGBARKEIT DER KREISTHEILUNGSGLEICHUNG.}
\noindent wo $\lambda,\mu',\nu'$ etc. Divisoren von $a^{\alpha-1}(a-1),\ldots$
\[\begin{aligned}
&\frac{\mu a^{\alpha-1}(a-1)}\lambda\text{ div. per }\mu'\\[5pt]
&\frac{\nu a^{\alpha-1}(a-1)}\lambda,\quad\frac{\nu\nu'{a'}^{\alpha'-1}(a'-1)}{\mu'}\text{ div. per }\nu''\text{ etc.}
\end{aligned}\]
\noindent [und] wo für
\[\begin{array}{rl}
x&\text{alle ganze Zahlen }0,1,2,\ldots,(a^{\alpha-1}(a-1):\lambda)-1\\[5pt]
x'&\text{alle ganze Zahlen }0,1,2,\ldots,({a'}^{\alpha'-1}(a'-1):\mu')-1\\[5pt]
x''&\text{alle ganze Zahlen }0,1,2,\ldots,({a''}^{\alpha''-1}(a''-1):\nu'')-1
\end{array}\]
\noindent zu substituiren sind. Hier besteht das ganze System aus
\[\frac{n(a-1)(a'-1)(a''-1)\ldots}{a\mathpunct{.}a'\mathpunct{.}a''\ldots\lambda\mathpunct{.}\mu'\mathpunct{.}\nu''\ldots}\]
\noindent Wurzeln und alle Wurzeln der Gleichung $X=0$ zerfallen in $\lambda\mu'\nu''$ etc.
Systeme, $\mu<\mu',\ \nu,\nu'<\nu''$ etc.
\begin{center}\textbf{BEMERKUNG.}\end{center}
\begingroup\small
\noindent\hspace*{\parindent}Nach der Tagebuchnotiz Nr. 136 hat \textsc{Gauss} die Unzerlegbarkeit
der Gleichung, der die primitiven Lösungen von $x^n=1$ genügen, für ein zusammengesetztes $n$
am 12. Juni 1808 bewiesen. Die vorstehende Aufzeichnung stammt wahrscheinlich aus dieser Zeit,
da in demselben Handbuche astronomische Rechnungen vorangehen, die die Örter der kleinen Planeten
Juno und Ceres für 1808 bestimmen. Vergl. auch die bezüglichen Bemerkungen in dem Artikel 14
des \textsc{Bachmann}schen Aufsatzes »Über \textsc{Gauss}’ zahlentheoretische Arbeiten«.
\par\noindent\hfill\textsc{Schlesinger}.\par
\endgroup
\ClosingRule
\clearpage
